% Électrolyse de l'eau : montage et bilan aux électrodes (L'essentiel, chapitre 10)
\begin{tikzpicture}[x=1cm,y=1cm]
\pic at (0,0) {becher={largeur=5.6,hauteur=2.7,niveau=2.25,liquide=pcEau}};
% Éprouvettes retournées, remplies de solution, avec un volume de gaz en haut
\foreach \x/\g in {-1.35/1.05,1.35/0.55}{%
  \fill[white] (\x-0.33,3.95-\g) rectangle (\x+0.33,3.95);
  \fill[pcEau] (\x-0.33,0.6) rectangle (\x+0.33,3.95-\g);
  \draw[pcSurface] (\x-0.33,3.95-\g) -- (\x+0.33,3.95-\g);
  \draw[pcVerre] (\x-0.39,0.55) -- (\x-0.39,3.89) arc[start angle=180,end angle=0,radius=0.39] -- (\x+0.39,0.55);}
% Électrodes
\filldraw[pcAppareil] (-1.47,0.35) rectangle (-1.23,2.0);
\filldraw[pcAppareil] (1.23,0.35) rectangle (1.47,2.0);
\draw[pcFil] (-1.35,2.0) -- (-1.35,2.52);
\draw[pcFil] (1.35,2.0) -- (1.35,2.52);
\draw[pcFil] (-1.35,0.35) -- (-1.35,-0.5) -- (-0.45,-0.5);
\draw[pcFil] (1.35,0.35) -- (1.35,-0.5) -- (0.45,-0.5);
\pcAppareilRond{(0,-0.5)}{G}
\node[font=\otSemi\fontsize{9}{10}\selectfont,text=otInk] at (-0.62,-0.82) {$-$};
\node[font=\otSemi\fontsize{9}{10}\selectfont,text=otInk] at (0.62,-0.82) {$+$};
\draw[pcElectron] (0.95,-0.92) -- (1.6,-0.92);
\node[pcLegende,anchor=west,text=otNavy] at (1.62,-0.92) {$\mathrm{e}^-$};
\draw[pcCourant] (-0.95,-0.14) -- (-1.6,-0.14);
\node[pcLegende,anchor=east,text=otRed!80!black] at (-1.62,-0.14) {$I$};
\pcNote[west]{(-1.1,3.6)}{(2.6,4.35)}{dihydrogène : \ce{2 H2O + 2 e- -> H2 + 2 HO-}}
\pcNote[west]{(1.62,3.5)}{(2.6,3.75)}{dioxygène : \ce{4 HO- -> O2 + 2 H2O + 4 e-}}
\node[pcLegendeG,anchor=east,align=right] at (-3.1,1.3) {cathode,\\reliée à la borne $-$ :\\réduction};
\node[pcLegendeG,anchor=west,align=left] at (3.1,1.3) {anode,\\reliée à la borne $+$ :\\oxydation};
\node[pcLegende,anchor=north] at (0,-1.15) {eau additionnée d'hydroxyde de sodium};
\end{tikzpicture}
\par\smallskip\legende{L'électrolyse de l'eau : le générateur force la transformation 2 H$_2$O $\rightarrow$ 2 H$_2$ + O$_2$ ; le volume de dihydrogène est double de celui de dioxygène.}
